% ---------------------------------------------------------------
% 2. Data Preprocessing and Cleaning
% Included from main.tex via % ---------------------------------------------------------------
% 2. Data Preprocessing and Cleaning
% Included from main.tex via % ---------------------------------------------------------------
% 2. Data Preprocessing and Cleaning
% Included from main.tex via % ---------------------------------------------------------------
% 2. Data Preprocessing and Cleaning
% Included from main.tex via \input{sections/preprocessing}
% ---------------------------------------------------------------

\subsection{Cleaning Rules}

The modelling datasets contain one row per delivered order (96,454 orders).
\todo{List each cleaning rule and how many rows it removes, e.g.\ orders that
were not delivered, missing delivery dates, invalid timestamps.}

\begin{table}[H]
    \centering
    \caption{Cleaning steps and rows retained.}
    \label{tab:cleaning}
    \begin{tabularx}{\textwidth}{|X|r|r|}
        \hline
        \textbf{Step} & \textbf{Rows removed} & \textbf{Rows remaining} \\
        \hline
        Raw orders & -- & 99,441 \\
        \hline
        \todo{step} & \tbd & \tbd \\
        \hline
        Final modelling dataset & -- & 96,454 \\
        \hline
    \end{tabularx}
\end{table}

\subsection{Missing Values and Outliers}

Zip-code coordinates are taken as the median latitude and longitude of each zip
prefix after removing geolocation points that fall outside Brazil.
\todo{Report the remaining missing values (e.g.\ 477 order distances, 265
customer coordinates) and how they are imputed inside the pipeline; describe
outlier handling.}

\subsection{Train/Test Split and Leakage Controls}

The data is split by purchase date so that the model is always evaluated on
orders placed after those it was trained on: training up to 30 June 2018
(83,947 orders), validation in July 2018 (6,156) and test in August 2018
(6,351). \todo{Note the July late rate of 3.4\% against 7.1\% in training and
6.2\% in test, and how this affects tuning.}

Columns known only after purchase are excluded as inputs: the carrier hand-off
and customer delivery dates, \texttt{delivery\_days}, \texttt{delay\_days},
\texttt{delivery\_status}, \texttt{review\_score}, payment approval time and
the seller shipping limit date. Each target is also excluded as an input to
the other problem.

% ---------------------------------------------------------------

\subsection{Cleaning Rules}

The modelling datasets contain one row per delivered order (96,454 orders).
\todo{List each cleaning rule and how many rows it removes, e.g.\ orders that
were not delivered, missing delivery dates, invalid timestamps.}

\begin{table}[H]
    \centering
    \caption{Cleaning steps and rows retained.}
    \label{tab:cleaning}
    \begin{tabularx}{\textwidth}{|X|r|r|}
        \hline
        \textbf{Step} & \textbf{Rows removed} & \textbf{Rows remaining} \\
        \hline
        Raw orders & -- & 99,441 \\
        \hline
        \todo{step} & \tbd & \tbd \\
        \hline
        Final modelling dataset & -- & 96,454 \\
        \hline
    \end{tabularx}
\end{table}

\subsection{Missing Values and Outliers}

Zip-code coordinates are taken as the median latitude and longitude of each zip
prefix after removing geolocation points that fall outside Brazil.
\todo{Report the remaining missing values (e.g.\ 477 order distances, 265
customer coordinates) and how they are imputed inside the pipeline; describe
outlier handling.}

\subsection{Train/Test Split and Leakage Controls}

The data is split by purchase date so that the model is always evaluated on
orders placed after those it was trained on: training up to 30 June 2018
(83,947 orders), validation in July 2018 (6,156) and test in August 2018
(6,351). \todo{Note the July late rate of 3.4\% against 7.1\% in training and
6.2\% in test, and how this affects tuning.}

Columns known only after purchase are excluded as inputs: the carrier hand-off
and customer delivery dates, \texttt{delivery\_days}, \texttt{delay\_days},
\texttt{delivery\_status}, \texttt{review\_score}, payment approval time and
the seller shipping limit date. Each target is also excluded as an input to
the other problem.

% ---------------------------------------------------------------

\subsection{Cleaning Rules}

The modelling datasets contain one row per delivered order (96,454 orders).
\todo{List each cleaning rule and how many rows it removes, e.g.\ orders that
were not delivered, missing delivery dates, invalid timestamps.}

\begin{table}[H]
    \centering
    \caption{Cleaning steps and rows retained.}
    \label{tab:cleaning}
    \begin{tabularx}{\textwidth}{|X|r|r|}
        \hline
        \textbf{Step} & \textbf{Rows removed} & \textbf{Rows remaining} \\
        \hline
        Raw orders & -- & 99,441 \\
        \hline
        \todo{step} & \tbd & \tbd \\
        \hline
        Final modelling dataset & -- & 96,454 \\
        \hline
    \end{tabularx}
\end{table}

\subsection{Missing Values and Outliers}

Zip-code coordinates are taken as the median latitude and longitude of each zip
prefix after removing geolocation points that fall outside Brazil.
\todo{Report the remaining missing values (e.g.\ 477 order distances, 265
customer coordinates) and how they are imputed inside the pipeline; describe
outlier handling.}

\subsection{Train/Test Split and Leakage Controls}

The data is split by purchase date so that the model is always evaluated on
orders placed after those it was trained on: training up to 30 June 2018
(83,947 orders), validation in July 2018 (6,156) and test in August 2018
(6,351). \todo{Note the July late rate of 3.4\% against 7.1\% in training and
6.2\% in test, and how this affects tuning.}

Columns known only after purchase are excluded as inputs: the carrier hand-off
and customer delivery dates, \texttt{delivery\_days}, \texttt{delay\_days},
\texttt{delivery\_status}, \texttt{review\_score}, payment approval time and
the seller shipping limit date. Each target is also excluded as an input to
the other problem.

% ---------------------------------------------------------------

\subsection{Cleaning Rules}

The modelling datasets contain one row per delivered order (96,454 orders).
\todo{List each cleaning rule and how many rows it removes, e.g.\ orders that
were not delivered, missing delivery dates, invalid timestamps.}

\begin{table}[H]
    \centering
    \caption{Cleaning steps and rows retained.}
    \label{tab:cleaning}
    \begin{tabularx}{\textwidth}{|X|r|r|}
        \hline
        \textbf{Step} & \textbf{Rows removed} & \textbf{Rows remaining} \\
        \hline
        Raw orders & -- & 99,441 \\
        \hline
        \todo{step} & \tbd & \tbd \\
        \hline
        Final modelling dataset & -- & 96,454 \\
        \hline
    \end{tabularx}
\end{table}

\subsection{Missing Values and Outliers}

Zip-code coordinates are taken as the median latitude and longitude of each zip
prefix after removing geolocation points that fall outside Brazil.
\todo{Report the remaining missing values (e.g.\ 477 order distances, 265
customer coordinates) and how they are imputed inside the pipeline; describe
outlier handling.}

\subsection{Train/Test Split and Leakage Controls}

The data is split by purchase date so that the model is always evaluated on
orders placed after those it was trained on: training up to 30 June 2018
(83,947 orders), validation in July 2018 (6,156) and test in August 2018
(6,351). \todo{Note the July late rate of 3.4\% against 7.1\% in training and
6.2\% in test, and how this affects tuning.}

Columns known only after purchase are excluded as inputs: the carrier hand-off
and customer delivery dates, \texttt{delivery\_days}, \texttt{delay\_days},
\texttt{delivery\_status}, \texttt{review\_score}, payment approval time and
the seller shipping limit date. Each target is also excluded as an input to
the other problem.
